\subsubsection{\'{E}tude d'un système ternaire Li-Cu-Cl}

\setcounter{equation}{0}
\setcounter{Q}{0}

On se propose d'étudier de manière thermodynamique le ternaire Li-Cu-Cl. 
\begin{center}
\begin{minipage}{0.48\textwidth}
\includegraphics[width=\textwidth]{Ex28/Fig01}
\end{minipage}\hfill\begin{minipage}{0.48\textwidth}
\captionof{figure}{Diagramme ternaire isotherme à \SI{25}{\celsius} où les corps purs, simples et complexes, sont indiqués}
\label{fig:28-01}
\end{minipage}
\end{center}

Sur la Fig. \ref{fig:28-01}, les lignes en trait plein indiquent les zones de démarcation entre des régions de composition différentes en chacune des phases. On cherchera à déterminer la f.e.m. où le système ternaire Li-Cu-Cl est au pôle $\oplus$ et la référence lithium métallique \ce{Li (s) | Li^+} à l'interface $\ominus$. Le système subit une évolution isotherme et isobare.

\Q Indiquer la composition dans toutes les régions du diagramme ternaire. Calculer la variance particulière et conclure.
\solution{On fait bien attention à indiquer la composition aussi sur les pôles et lignes de démarcation, sinon la réponse est incomplète. On peut alors écrire : 
\begin{center}
\includegraphics[width=0.8\textwidth]{Ex28/Corr03}
\end{center}
Tous les domaines où $v_p = 0$ auront une f.e.m. constante (aucun degré de liberté intensif). Les sauts de potentiel auront lieu sur les lignes de démarcation où $v_p = 1$.}

\begin{minipage}{0.48\textwidth}
\includegraphics[width=\textwidth]{Ex28/Fig02}
\end{minipage}\hfill\begin{minipage}{0.48\textwidth}
\captionof{figure}{Ligne de composition lorsque la fraction de lithium varie au sein d'une pile \ce{Li||CuCl}}
\end{minipage}

On considère dans un second temps que l'on étudie la pile suivante : 
\begin{align*}
\ominus ~ \ce{Li} | \ce{LiCl} | \ce{CuCl} ~ \oplus
\end{align*}

\Q Où se trouve-t-on dans le diagramme lorsque la pile est chargée et lorsqu'elle est déchargée ?
\solution{Au début, lorsque la pile est chargée, on se trouve au pôle \ce{CuCl}. Ensuite, en déchargeant progressivement, on libère des ions lithium vers le pôle $\oplus$. On parcourt la ligne en pointillé vers le pôle chimique $\ce{Li}$ (qui n'est jamais atteint physiquement).}

\Q Représenter le profil de potentiel associé à la décharge pour une st{\oe}chiométrie de lithium comprise dans l'intervalle $[0,1]$.
\solution{On utilise la relation entre enthalpie et travail chimique, qui nous permet d'écrire : 
\begin{align*}
\Delta_r G = - n F e = -n F E_\oplus [\si{\volt\per{\ce{Li^+} / \ce{Li}}}]
\end{align*}
On sait que dans le domaine B la variance est nulle. On peut donc calculer la valeur du plateau de f.e.m. dans ce domaine. Commençons par modéliser l'équation de décharge : 
\begin{center}
\begin{minipage}{0.4\textwidth}
\begin{align*}
\ce{Li &= Li^+ + e^-} \\
\ce{CuCl + e^- &= CuCl + Cl^-} \\
\hline
\ce{CuCl + Li &= Cu + Li^+ + Cl^-} \\
\ce{CuCl + Li &= Cu + LiCl}
\end{align*}
\end{minipage}
\end{center}
On peut calculer les enthalpies de formation de Gibbs à \SI{298}{\kelvin} du domaine B : 
\begin{align*}
\Delta_r G^\circ ( \text{B} ) &=  \Delta_f G^\circ ( \ce{LiCl} ) - \Delta_f G^\circ ( \ce{CuCl} ) \\
& \stackrel{AN}{=} -245.3 ~ \si{\kilo\joule\per\mole}
\end{align*}
On peut ainsi écrire le potentiel cathodique du domaine B comme :
\begin{align*}
E_\oplus ( \text{B} ) &= -\frac{\Delta_r G^\circ ( \text{B} ) }{F} \stackrel{AN}{=} 2.542 ~ \si{\volt\per{\ce{Li^+ / Li}}}
\end{align*}
On s'attend à un profil comme suit :
\begin{center}
\includegraphics[width=0.2\textwidth]{Ex28/Corr01}
\end{center}}


\begin{minipage}{0.48\textwidth}
\includegraphics[width=\textwidth]{Ex28/Fig03}
\end{minipage}\hfill\begin{minipage}{0.48\textwidth}
\captionof{figure}{Ligne de composition lorsque la fraction de lithium varie au sein d'une pile \ce{Li||CuCl2}}
\end{minipage}


On considère dans un second temps que l'on étudie la pile suivante : 
\begin{align*}
\ominus ~  \ce{Li} | \ce{LiCl} | \ce{CuCl2} ~ \oplus
\end{align*}

\Q Où se trouve-t-on dans le diagramme lorsque la pile est chargée et lorsqu'elle est déchargée ?
\solution{Au début, lorsque la pile est chargée, on se trouve au pôle \ce{CuCl2}. Ensuite, en déchargeant progressivement, on libère des ions lithium vers le pôle $\oplus$. On parcourt la ligne en pointillé vers le pôle chimique $\ce{Li}$ (qui n'est jamais atteint physiquement).}

\Q Représenter le profil de potentiel associé à la décharge pour une st{\oe}chiométrie de lithium comprise dans l'intervalle $[0,2]$.
\solution{On utilise la relation entre enthalpie et travail chimique, qui nous permet d'écrire : 
\begin{align*}
\Delta_r G = - n F e = -n F E_\oplus [\si{\volt\per{\ce{Li^+} / \ce{Li}}}]
\end{align*}
On sait que dans les domaines C et B la variance est nulle. On peut donc calculer la valeur du plateau de f.e.m. dans ces domaines. Commençons par modéliser les équations de décharge : 
\begin{itemize}
\item dans le domaine C : 
\begin{center}
\begin{minipage}{0.4\textwidth}
\begin{align*}
\ce{Li &= Li^+ + e^-} \\
\ce{CuCl2 + e^- &= CuCl + Cl^-} \\
\hline
\ce{CuCl2 + Li &= CuCl + Li^+ + Cl^-} \\
\ce{CuCl2 + Li &= CuCl + LiCl}
\end{align*}
\end{minipage}
\end{center}
\item dans le domaine B : 
\begin{center}
\begin{minipage}{0.4\textwidth}
\begin{align*}
\ce{Li &= Li^+ + e^-} \\
\ce{CuCl + e^- &= Cu + Cl^-} \\
\hline
\ce{CuCl + Li &= Cu + Li^+ + Cl^-} \\
\ce{CuCl + Li &= Cu + LiCl}
\end{align*}
\end{minipage}
\end{center}
\end{itemize}
On peut alors calculer les enthalpies de formation de Gibbs à \SI{298}{\kelvin} dans chacun des domaines : 
\begin{itemize}
\item dans le domaine C :  
\begin{align*}
\Delta_r G^\circ ( \text{C} ) &=  \Delta_f G^\circ ( \ce{LiCl} ) + \Delta_f G^\circ ( \ce{CuCl} ) - \Delta_f G^\circ ( \ce{CuCl2}  ) \\
&\stackrel{AN}{=} -348.9 ~ \si{\kilo\joule\per\mole}
\end{align*}
\item dans le domaine B : 
\begin{align*}
\Delta_r G^\circ ( \text{B} ) &=  \Delta_f G^\circ ( \ce{LiCl} ) - \Delta_f G^\circ ( \ce{CuCl} ) \\
& \stackrel{AN}{=} -245.3 ~ \si{\kilo\joule\per\mole}
\end{align*}
\end{itemize}
On peut alors écrire le potentiel cathodique dans les domaines B/C comme :
\begin{align*}
E_\oplus ( \text{C} ) &= -\frac{\Delta_r G^\circ ( \text{C} ) }{F} \stackrel{AN}{=} 3.615 ~ \si{\volt\per{\ce{Li^+ / Li}}} \\
E_\oplus ( \text{B} ) &= -\frac{\Delta_r G^\circ ( \text{B} ) }{F} \stackrel{AN}{=} 2.542 ~ \si{\volt\per{\ce{Li^+ / Li}}}
\end{align*}
On s'attend à un profil comme suit :
\begin{center}
\includegraphics[width=0.4\textwidth]{Ex28/Corr02}
\end{center}}

\paragraph*{Données}
\begin{itemize}
\item Enthalpie libre de formation à \SI{298}{\kelvin}
\begin{center}
\begin{tabular}{cc}
\toprule
Phase & $\Delta_f G^\circ$ \\
& [\si{\kilo\joule\per\mole}] \\
\midrule
\ce{LiCl} & -384.0 \\
\ce{CuCl} & -138.7 \\
\ce{CuCl2} & -173.8 \\
\bottomrule
\end{tabular}
\end{center}
\end{itemize}